\chapter{Contribuição esperada e limites}\label{produto-publicuxe1vel-e-limites-de-escopo}

\textbf{Título provisório do artigo:} \emph{Frequency compression and
prior sensitivity in pulsar-timing tests of massive gravitational
waves}.

Uma contribuição publicável poderá consistir em:

\begin{itemize}
\tightlist
\item
  critérios quantitativos para empregar uma frequência efetiva sem
  alterar materialmente a inferência;
\item
  medição de viés, perda de informação ou falsa preferência por um
  modelo em cenários fisicamente especificados;
\item
  demonstração de como uma componente escalar vinculada modifica esses
  resultados;
\item
  benchmarks abertos e validação independente que permitam reproduzir a
  conclusão.
\end{itemize}

O manuscrito deve distinguir claramente resultados novos de reproduções.
A aceitação dependerá da originalidade efetivamente confirmada, da
qualidade da validação e da relevância dos efeitos encontrados.
\emph{Classical and Quantum Gravity} e \emph{Physical Review D} são
alvos temáticos plausíveis, como mostram os antecedentes citados; a
escolha final depende do resultado e da orientação.

Se a compressão se mostrar adequada, o artigo poderá apresentar seu
domínio de validade e limites superiores para o viés. Se a revisão
encontrar trabalho idêntico, o recorte deverá migrar cedo para a
extensão escalar ou para outro mecanismo concreto de perda de
informação. O projeto não deve se expandir simultaneamente para
cosmologia completa, emissão não linear de binárias, LISA, SKA e
bispectro.

\textbf{Recomendação:} preservar do TG a ligação entre polarizações,
massa e baixas frequências, e transformar essa base em um mestrado de
validação teórica e numérica dos testes atuais com PTAs. O diferencial
candidato está na confiabilidade e na informação das inferências, e não
na redescoberta de modos de polarização já conhecidos.
